\begin{center}
    \section*{\fontsize{20}{20}\selectfont Chapter 2}
\end{center}
\vspace{10mm}
\section{Literature Review}

\subsection{Introduction}

This chapter surveys the four research threads that the framework in this thesis brings together: vision-based fire and smoke detection, acoustic fire suppression, delay-tolerant and mesh-based messaging for alerting, and computer-aided dispatch for fire and emergency services. A fifth, shorter subsection considers the Bangladesh-specific fire-response literature, which is sparse but relevant. The chapter closes by identifying the gap that the framework is designed to fill.

\subsection{Vision-Based Fire and Smoke Detection}

Vision-based fire detection has evolved through three broad phases. The broader context for this evolution is the long-running shift in video surveillance from purely human-monitored, recording-only systems toward architectures that incorporate automated visual analysis, a shift surveyed by Tsakanikas and Dagiuklas \cite{Tsakanikas2018} and by Shah, Javed, and Shafique \cite{Shah2007}, who noted that most installations remain fundamentally passive, serving as forensic tools rather than active safety mechanisms. The \textbf{earliest phase} used rule-based pixel methods, in which a frame was classified as containing fire if a sufficient fraction of its pixels fell within a particular range of red, orange, and yellow hues, with high saturation and high value in the HSV colour space. Celik and Demirel \cite{celik2010fire} provided a representative and widely-cited formulation of this approach, and the heuristic fallback in our detection layer is a direct descendant of their rule set, extended with connected-component labelling to produce bounding boxes rather than a single binary decision.

The \textbf{second phase} introduced hand-crafted features combined with classical machine-learning classifiers. Color histograms, texture descriptors, and motion vectors extracted from consecutive frames were fed into support vector machines or random forests. Rossi and Acqua \cite{rossi2015feature} surveyed this body of work and noted that the approaches gained roughly ten percentage points of accuracy over the rule-based methods but at a substantial computational cost, which limited their deployment on edge devices.

The \textbf{third phase}, which dominates the current literature, is deep learning. Convolutional neural networks have been applied to fire detection since at least 2016. The VGG and ResNet families were used initially as image classifiers, with the network producing a single fire-or-no-fire label for the whole frame \cite{simonyan2014vgg}. The shift to object-detection architectures, which produce bounding boxes around fire and smoke regions, came a few years later. The YOLO family \cite{ultralytics2023yolov8}, the Single Shot Detector family, and Faster R-CNN have all been applied. \textbf{YOLOv8} has become the dominant choice for resource-constrained deployment because of its favourable accuracy-to-latency trade-off at the nano and small scales. Yasar and Yildiz \cite{yasar2024fireyolo} reported mAP@0.5 of 0.93 on a combined fire and smoke benchmark with YOLOv8n at 30 frames per second on a Jetson Nano, and the D-Fire benchmark \cite{dunnhofer2023dfire} has become the de facto evaluation set for the task. A related line of work by Ahmed and Echi \cite{Ahmed2021} presented Hawk-Eye, an AI-powered threat detector for intelligent surveillance cameras, and Poirier \cite{poirier2024anomaly} examined real-time anomaly detection in video streams more broadly, highlighting the difficulty of generalising anomaly models across varied, uncontrolled environments.

A separate line of work has explored transformer-based detectors for fire, on the grounds that the global attention mechanism may capture the diffuse and irregular shape of smoke better than the local receptive field of a CNN. Liang et al. \cite{liang2024firetransformer} reported a modest improvement over YOLOv8 at the cost of a 4$\times$ increase in inference time. For the deployment scenario targeted in this thesis, in which the detector runs in a browser through onnxruntime-web, the inference-time penalty is unacceptable, and we did not pursue the transformer line further.

Edge deployment has been a recurring concern. Park et al. \cite{park2017scalable} described a scalable architecture for surveillance using edge computing, in which inference was pushed to camera-attached devices rather than to a central server. Dautov et al. \cite{dautov2018metropolitan} generalised this approach to metropolitan-scale IoT surveillance. The framework in this thesis follows the same edge-first principle, but it adds the network-independent relay, the acoustic suppressor, and the dispatch workflow that the edge-surveillance literature does not address.

A smaller body of work has considered the ONNX runtime as a deployment target for fire detectors. Mahmoud and Abouelnaga \cite{mahmoud2023onnx} benchmarked onnxruntime-web against TensorFlow.js for several YOLO variants and reported a 1.5 to 2 times speed-up on Chromium-based browsers, with a smaller but still significant speed-up on Firefox. Their finding informed our choice of onnxruntime-web as the inference engine for the detection layer.

\subsection{Acoustic Fire Suppression}

The interaction between sound fields and flames is a much smaller literature than the fire-detection literature, but it is the literature most directly relevant to the suppression layer of this thesis. The phenomenon itself is old: Lord Rayleigh noted in \textit{The Theory of Sound} that a flame could be made to flicker, lengthen, or extinguish in response to an acoustic field, and the 19th-century demonstration of a ``singing flame'' is a staple of physics lectures.

The modern literature begins with the \textbf{2012 DARPA demonstration} in which a pair of large loudspeakers emitting tones in the 30 to 60 Hz band extinguished a small flame at close range \cite{darpa2012acoustic}. The DARPA work was widely reported in the popular press but treated with scepticism in the combustion-physics community, partly because the demonstration conditions were not fully documented and partly because the mechanism was not immediately obvious. Subsequent work by \textbf{Kim and Park} \cite{kim2015acoustic} at the University of Georgia clarified the mechanism: the dominant effect is \textbf{velocity coupling at the flame base}. The oscillating flow thins the flame front below the quenching thickness, while the periodic pressure variation disrupts the heat-feedback loop between the flame and the fuel surface. The combined effect is that the flame's local burning rate drops below the sustenance threshold, and the flame goes out.

McKinney and co-workers \cite{mckinney2019sonic} extended Kim and Park's work with a parametric study of frequency, intensity, and wave type. They reported that the \textbf{30 to 60 Hz band is the most effective} across a range of flame sizes, that sine waves outperform square and sawtooth waves at the same intensity (because the latter's high-frequency components are absorbed by the air before reaching the flame), and that the intensity threshold for extinction scales with the cube of the flame's characteristic length. The cube scaling implies that the method is most effective on small flames and rapidly becomes impractical for room-scale fires, which is consistent with the DARPA demonstration's small-flame focus.

More recent work has considered practical implementations. \textbf{Tran and Liao} \cite{tran2021portable} built a portable acoustic extinguisher driven by a 200-watt class-D amplifier and a neodymium-driver horn loudspeaker, and reported extinction of a 15-cm alcohol pan fire in 3.8 seconds at 45 Hz. Their hardware is the closest published precedent to the suppression layer in this thesis, although they did not integrate it with a detection or dispatch pipeline. \textbf{Li et al.} \cite{li2022drone} mounted an acoustic extinguisher on a quadcopter and demonstrated remote suppression of a small flame, again without integration with a detection pipeline.

A separate line of work has considered the safety of acoustic suppression. The 30 to 60 Hz band is below the human pain threshold (which begins at roughly 120 dB at 1000 Hz and rises at lower frequencies), but sustained exposure above 100 dB can cause hearing damage. The suppression layer in this thesis is therefore designed to fire only on a confirmed detection, with a maximum duration of 10 seconds, and with a safety interlock that requires the operator to have acknowledged the alert before the suppressor can fire in occupied spaces.

To our knowledge, no prior work has integrated an acoustic fire suppressor into an autonomous fire-detection and dispatch pipeline. The contribution of this thesis on the suppression front is therefore the integration, not the suppressor itself.

\subsection{Reactive and Proactive Detection Systems}

A distinction that recurs in the surveillance literature is between systems that merely detect and log an event and systems that actively trigger a response. Jung and Kim \cite{jung2025opencv} analysed the security vulnerabilities of OpenCV-based YOLOv10 IP camera systems for disaster safety management and noted that such systems depend on stable IP connectivity, which fails precisely when an emergency is in progress. Katragada et al. \cite{katragada2023atm} applied a multimodel YOLO architecture to real-time ATM surveillance with near 99 percent class-specific detection accuracy, but their system functions as a notification tool and does not incorporate an automated escalation mechanism. The pattern, high detection accuracy paired with a single internet-dependent notification path, is consistent across much of the reviewed literature and represents a structural limitation rather than a deficiency of any individual model.

A smaller body of work has tackled the proactive side. Tauseef et al. \cite{tauseef2024smartguard} described SmartGuard, a YOLOv8 and SORT system for unattended-object detection in transport hubs, which sends push notifications to a configured phone list but does not address network failure. Behera et al. \cite{behera2025lstm} integrated YOLOv8 with an LSTM for temporal pattern recognition, with a reported mAP@0.5 of 0.90, but again the alerting path is a single push notification. The Amanot-Net project \cite{amanotnet2026}, on which this thesis builds, added a dual-channel alert mechanism (SIP plus LoRaWAN) to a YOLOv8-based weapon detector, but it stopped at the alert and did not enter a dispatch workflow or attempt on-site suppression. The framework in this thesis extends Amanot-Net in three ways: it replaces the LoRaWAN backup with a peer-to-peer mesh that reuses the local Wi-Fi spectrum, it adds the acoustic-wave suppressor that acts on the fire while the alert is in flight, and it adds the dispatch workflow on the back end.

\subsection{Delay-Tolerant Networking and Mesh Messaging}

The alerting layer of the framework draws on the delay-tolerant networking and mesh-messaging literature. The DTN line begins with \textbf{Vahdat and Becker's epidemic routing} \cite{vahdat2000}, which floods messages whenever two nodes meet and which is asymptotically optimal for delivery ratio at the cost of quadratic resource usage. Spyropoulos et al. \cite{spyropoulos2005} proposed Spray-and-Wait, which bounds the number of copies of a message to a fixed $L$ and improves on epidemic's resource profile. The Bundle Protocol was standardised at the IETF as RFC 5050 in 2007 and revised as RFC 9171 in 2022 \cite{rfc5050,rfc9171}. These works establish the store-and-forward, hop-count-limited flooding, and UUID-based deduplication primitives that our mesh relay uses.

A parallel line of work considered practical mesh messaging on commodity devices. The \textbf{Serval project} \cite{serval2013} demonstrated mesh telephony on Android phones with custom firmware, using Wi-Fi ad-hoc mode and a peer-to-peer voice protocol. \textbf{Briar} \cite{briar2018} provided peer-to-peer messaging over Tor hidden services and Bluetooth. \textbf{Firechat} \cite{firechat2015} briefly achieved fame during the 2014 Hong Kong protests by using Apple's Multipeer Connectivity framework, but it was discontinued in 2015. \textbf{Manyverse} \cite{manyverse2019} brought Secure Scuttlebutt to Android for offline social networking. None of these systems run on OpenWrt, and all require the user to install a custom application before the network can be used.

The \textbf{Mesh-Chat project} \cite{meshchat2026} closed this gap by demonstrating an offline peer-to-peer messenger that runs on a stock OpenWrt router, occupies 19 kilobytes of package space and 1.7 megabytes of resident memory, and serves a web UI from the router itself so that no client application is required. Mesh-Chat evaluated four message-propagation strategies (flood, gossip, hybrid, mesh-aware) on a 200-node discrete-event simulator and showed that gossip reduces amplification by an order of magnitude at large $N$ with a modest latency cost. The mesh relay in this thesis is a direct extension of Mesh-Chat, modified to flood compact JSON incident alerts rather than free-text chat messages, and to use Ed25519 signing for per-router authentication.

A separate line of work has considered LoRaWAN as a backup channel for emergency alerting. Sartori et al. \cite{sartori2024lora} described a LoRaWAN-based fire-alarm system for remote buildings, in which a smoke sensor triggered a LoRaWAN uplink to a network server, which then forwarded an alert to a phone app. The approach is sound but depends on LoRaWAN gateway coverage, which is uneven in Bangladesh. Wi-Fi mesh, by contrast, runs on hardware that the building already owns, which motivated our choice of Wi-Fi over LoRaWAN for the relay layer.

\subsection{Computer-Aided Dispatch for Fire Services}

Computer-aided dispatch has been in continuous use in public-safety answering points since the 1980s. The \textbf{Next Generation 911 program} in the United States \cite{ng9112023} has spent more than a decade migrating the original circuit-switched 911 system to an IP-based, geospatially aware platform, in which the caller's location is conveyed as a structured data element rather than read out over the phone. The NG 911 data model, documented by the National Emergency Number Association, includes an incident record, a resource record, and an event log that together capture the lifecycle of an emergency from intake to closure.

The European equivalent is the \textbf{eCall system}, mandatory in new vehicles in the European Union since 2018, which automatically places a 112 call to the nearest public-safety answering point with a structured data payload containing the vehicle's location and crash details \cite{ecall2018}. The eCall data model and the NG 911 data model are broadly compatible, and both informed the incident schema used in this thesis.

The South Asian picture is more fragmented. \textbf{India's 112 emergency number}, launched in 2017, consolidates police, fire, and ambulance calls at state-level command centres, but vehicle tracking and inter-agency escalation remain works in progress \cite{india1122022}. \textbf{Bangladesh's 999 service}, launched in the same year, handles call-taking well but does not yet offer automatic dispatch, live vehicle tracking, or structured inter-agency coordination \cite{bangladesh9992018}. The Fire Service and Civil Defence's own modernization program, set out in its Fire Service and Civil Defence Strategic Plan 2021 to 2025 \cite{fscd2021strategic}, identifies dispatch automation, station directory modernisation, and a citizen-facing incident-reporting channel as priority areas. The dispatch layer of the framework in this thesis is designed against these priorities.

A small academic literature has considered dispatch-system design for low-resource settings. Khan et al. \cite{khan2022cad} described a prototype CAD system for the Lahore rescue service that handled automatic vehicle assignment by geo-distance, with a reported 35 percent reduction in dispatch time over the manual baseline. Ali and Hossain \cite{ali2023dispatch} simulated a similar system for Dhaka and reported a comparable reduction. Neither system, however, was integrated with a detection layer, a suppression layer, or a network-independent relay, which is the gap this thesis addresses.

\subsection{IoT-Based Fire Sensor Networks}

A separate strand of literature has considered IoT-based fire sensor networks, in which low-cost smoke, temperature, humidity, and gas sensors are connected through a wireless protocol such as Wi-Fi, Zigbee, or LoRaWAN to a backend that aggregates readings and triggers alarms. Sinha et al. \cite{sinha2022iotfire} surveyed this body of work and identified two dominant architectures: a \textbf{star topology} with a single gateway, and a \textbf{mesh topology} with multiple relaying nodes. The star topology is simpler but has a single point of failure at the gateway, while the mesh topology is more robust but more complex to configure. A related body of work has considered the broader cybersecurity and access-control dimensions of IoT-based safety systems. Zhukabayeva et al. \cite{zhukabayeva2025edge} proposed an edge-computing-based framework for network traffic analysis and intrusion detection in industrial IoT, illustrating that edge-level processing can meaningfully reduce reliance on centralized infrastructure. Shahzaib, Ahmed, and Sharif \cite{shahzaib2025doorlock} presented a biometric and GSM-enabled smart door lock system that combines physical access control with cellular-network alerting, although such systems are typically event-triggered and rely on a single communication channel. Razzaq and Shah \cite{razzaq2025cybersecurity} conducted a scientometric analysis of machine-learning-driven cybersecurity research, observing a broader trend toward decentralized, intelligent defense mechanisms that is relevant to the resilience of any IoT-based fire-response pipeline.

LoRaWAN has been a particularly popular choice for fire sensor networks in rural and peri-urban deployments, on account of its long range and low power consumption. Cao et al. \cite{cao2023lorafire} described a LoRaWAN-based forest-fire detection system in China that achieved a 12-kilometre range with a 250 milliampere battery. The system detected fires reliably but produced alerts only at the gateway, with no end-to-end dispatch integration. The framework in this thesis uses IoT-style environmental sensors as an auxiliary input to the detection layer, but the primary detection signal is the camera feed, which provides richer information than a smoke sensor alone.

NB-IoT has been considered as an alternative to LoRaWAN, particularly in urban deployments where cellular coverage is good. Wang et al. \cite{wang2024nbiot} compared NB-IoT and LoRaWAN for smart-city fire detection and found that NB-IoT had lower latency in dense urban areas but consumed more power and required a subscription to a cellular network. For the deployment scenario targeted in this thesis, in which the building's internet connection may be down, NB-IoT shares the same failure mode as the primary channel and was therefore not pursued.

\subsection{The Bangladesh Fire-Response Context}

The Bangladesh-specific literature is sparse but operational. The Bangladesh Fire Service and Civil Defence's annual report \cite{fscdAnnual2024} records roughly 20{,}000 fire incidents per year nationally, with approximately 200 to 300 fatalities. The Dhaka metropolitan area accounts for roughly 40 percent of the incidents and 50 percent of the fatalities, reflecting the higher density and the older building stock in the capital. The report identifies three structural problems: \textbf{delayed reporting}, \textbf{inadequate station coverage} in the older wards, and the \textbf{absence of an automatic dispatch system}.

Islam and Hossain \cite{islam2020dhakafire} analysed the response-time distribution for Dhaka and reported a \textbf{median of 14 minutes} from ignition to first appliance on scene, against an international benchmark of 5 to 8 minutes. Their analysis identified the call-taking and station-assignment phases as the dominant contributors to the delay, with travel time a secondary factor. This finding is the operational basis for the dispatch layer's automatic station-assignment feature.

The Fire Service and Civil Defence's strategic plan \cite{fscd2021strategic} sets out a five-year modernisation agenda that includes the deployment of a computer-aided dispatch system, the digitisation of the station directory, the introduction of a citizen-facing incident-reporting channel, and the integration of GPS-based vehicle tracking. The framework in this thesis is consistent with this agenda and is intended as a working prototype that the Fire Service can evaluate and extend.

\subsection{Gaps in the Existing Research}

Synthesising the reviewed literature reveals a consistent pattern. The fire-detection thread has matured to the point where edge-deployable YOLOv8 models reach better than 0.90 mAP at interactive frame rates. The acoustic-suppression thread has demonstrated that low-frequency sound fields can extinguish small flames in laboratory conditions. The mesh-messaging thread has produced deployable systems on OpenWrt that survive the failure of the public internet. The computer-aided dispatch thread has standardised incident schemas and operational workflows in the NG 911 and eCall models. The Bangladesh-specific thread has identified the operational gaps in the current Fire Service practice.

What is missing is the integration. No reviewed system combines a YOLOv8-based edge detector with an acoustic-wave suppressor, a mesh-based alert relay, and a structured dispatch workflow in a single pipeline with a shared incident schema. The detection systems stop at the alert. The suppression systems are laboratory demonstrations without detection or dispatch. The mesh systems do not consider fire response. The dispatch systems assume that the alert has already arrived. The framework in this thesis is designed to fill that gap, by extending the Amanot-Net edge-first dual-channel pattern with a suppression layer and a dispatch layer, and by binding the four layers with a single incident schema.

\section{Problem Analysis}

Building on the literature reviewed above, the core problem addressed by this research can be decomposed into five interrelated issues.

\begin{enumerate}
\item \textbf{The detection-action disconnect.} Across the reviewed systems \cite{thakur2024weapon,katragada2023atm,behera2025lstm,tauseef2024smartguard,chung2024yololsd}, AI-based detection is frequently treated as the end goal rather than the starting point of a response pipeline. Even highly accurate detectors function largely as notification tools, leaving the actual escalation and human response process manual or only loosely automated. This results in a high Mean Time To Response, since a detected threat does not automatically translate into a timely, dependable alert reaching the appropriate personnel.

\item \textbf{Single point of failure in communication.} Both vision-based systems and IoT-based access-control systems are typically designed around a single communication channel, most commonly a standard IP network or a GSM link. As Sterbenz et al. \cite{sterbenz2010resilience} emphasise, any system lacking redundancy in its communication path is inherently vulnerable to a single point of failure: if that one channel is congested, jammed, or down, the entire alerting function collapses regardless of how accurate the underlying detection model is.

\item \textbf{Network dependency in infrastructure-constrained environments.} Field evaluations such as that of Yao et al. \cite{yao2025field} and industry assessments such as Verizon's safe-cities report both indicate that real-world deployment conditions, particularly intermittent connectivity, power instability, and bandwidth constraints, significantly undermine systems that assume continuous, high-bandwidth network access. This is especially relevant to unstructured public spaces in developing regions such as Bangladesh, where such conditions are the norm rather than the exception.

\item \textbf{No on-site suppression between detection and crew arrival.} Even when detection is fast and the alert is delivered, the fire continues to grow until a crew arrives. The acoustic-suppression literature \cite{darpa2012acoustic,kim2015acoustic,mckinney2019sonic,tran2021portable} demonstrates that a small flame can be extinguished in seconds by a low-frequency sound field, but no reviewed system integrates this capability into an autonomous pipeline. The window between detection and crew arrival is the window in which a small flame is most easily extinguished, and the framework in this thesis is designed to act in that window.

\item \textbf{Lack of integration with dispatch workflows.} The fire-detection literature rarely considers the workflow that follows the alert. The dispatch literature rarely considers the source of the alert. The result is that the two systems are designed in isolation, and the integration is left to ad-hoc engineering at deployment time. The framework in this thesis treats the integration as a first-class design problem and provides a shared incident schema that binds the two.
\end{enumerate}

Taken together, these five problems define the central research gap: the absence of an integrated framework that performs real-time, context-aware fire detection, attempts on-site acoustic suppression, relays the alert over a network-independent channel, and feeds the alert into a structured dispatch workflow in a single pipeline. The framework proposed in the next chapter is designed to address this gap directly.

\subsection{Summary}

This chapter has surveyed the four research threads that the framework brings together, has identified the gaps in the existing literature, and has decomposed the core problem into five interrelated issues. The next chapter presents the framework in detail.
